\section{Management Information System}
{\color{sub}\footnotesize 5 hrs \gap$\cdot$\gap asked in \mk{24 of 29 papers}
\gap$\cdot$\gap 8 syllabus topics \gap$\cdot$\gap \mk{248 marks, 10.5\% of the archive}
$\rightarrow$ about 6 of every 80 marks: one 8-mark question, split 3+5 or 4+4.}

\vspace{3pt}\noindent
\colorbox{gdL}{\makebox[\dimexpr\textwidth-2\fboxsep][l]{%
  \textcolor{gd}{\bfseries\footnotesize READ THIS FIRST}}}
\par\nobreak\vspace{3pt}

MIS questions almost always open with \mk{``define MIS / why is MIS needed?''} (17 papers)
and add one of three second halves:
\begin{enumerate}[label=\arabic*., leftmargin=1.4em, itemsep=1pt]
  \item \textbf{Types of information systems} (TPS, MIS, DSS, ESS) and DSS vs MIS
    (7 papers) $\rightarrow$ one table in $\S$5.3.
  \item \textbf{Information support for functional areas} (5 papers) or \textbf{at different
    levels} $\rightarrow$ one pyramid figure in $\S$5.3 answers both.
  \item \textbf{Planning and decision making} (7 papers), \textbf{computers and MIS},
    \textbf{database} and \textbf{networking} systems $\rightarrow$ $\S$5.2 and $\S$5.4.
\end{enumerate}
The 2065--2068 papers asked ``hierarchy of information needs'' and ``information
architecture''; only S.K.\ Joshi's textbook covers them, and they are summarized from it.

Tier chips count \textbf{papers out of 29}: \tS{n} 8+ \gap \tF{n} 4--7 \gap \tP{n} 1--3
inside an 8-mark question \gap \tL{n} 1--3, short note only \gap \tN{} never set.
Bold year = Regular exam.

\hr

%=======================================================================
\T{5.1 Data, Information and MIS}

\Q{Define MIS. Why is MIS needed? Its functions and importance}{\m{2}\m{3}\m{4}\m{5}}
{\color{sub}\footnotesize \tS{17} \yr{82 Ba, 81 Ba, \textbf{80 Bh}, 79 Ba, \textbf{76 Ch}, 76 Ash, 74 Ash, 73 Shr, \textbf{71 Ch}, \textbf{70 Ch}, 70 Asa, \textbf{69 Ch}, \textbf{68 Ch}, \textbf{68 Ba}, \textbf{67 Asa}, \textbf{66 Bh}, \textbf{65 Shr}} \gap$\cdot$\gap ``needs for MIS'' short note in 68 Ch, 66 Bh, 65 Shr; also the second half of 73 Ch and 71 Ch (filed with case studies)}

\sbs{%
\begin{itemize}
  \item \textbf{Management}: directing an enterprise toward objectives.
    \textbf{Information}: processed data that gives knowledge. \textbf{System}: interacting
    procedures forming an organized whole.
  \item \textbf{MIS}: ``a system of obtaining, abstracting, storing and analysing data to
    produce information for planning, controlling and decision making, given to managers at
    the time they can best use it.''
  \item Henry C.\ Lucas: ``a set of organized procedures which, when executed, provide
    information to support decision making.''
  \item Uses computing and communication technology; draws on internal and external
    sources; covers past, present and projected activity.
\end{itemize}
}{%
\lead{Functions}
\begin{enumerate}
  \item \textbf{Data capture}: from internal and external sources.
  \item \textbf{Processing}: classify, calculate, compare, summarize.
  \item \textbf{Storage and retrieval}.
  \item \textbf{Determining information needs}: what, how much, when, for whom.
  \item \textbf{Evaluation and abstraction}: judge reliability, edit, condense.
  \item \textbf{Dissemination}: right information, right manager, right time.
  \item \textbf{Prediction, planning, control} support.
\end{enumerate}
}

\sbs{%
\lead{Need for MIS}
\begin{itemize}
  \item Every manager decides on information; more complex firms $\rightarrow$ more data
    than a person can handle.
  \item \textbf{Production managers}: production, labour, machine, overhead costs;
    capacity.
  \item \textbf{Marketing managers}: sales trends, market research, selling costs, new
    products.
  \item \textbf{Personnel managers}: turnover, absenteeism, skills, wage levels, labour
    market.
  \item \textbf{Why MIS in addition to task software (79 Ba)}: payroll, billing, inventory
    packages each serve one task; MIS \textbf{integrates} their data into summary and
    exception reports across the firm, for decisions no single package supports.
\end{itemize}}{%
\lead{Importance}
\begin{itemize}
  \item The ``heart'' of the firm: timely, relevant information for \textbf{decisions}.
  \item \textbf{Measures performance against goals} $\rightarrow$ control.
  \item Supports \textbf{planning}, forecasting, coordination across departments.
  \item Reduces uncertainty and cost of information; keeps it current.
  \item Stores confidential information securely for future use.
  \item Competitive advantage: faster response to market.
  \item Eg a bank's MIS flags branches whose loan defaults exceed target; an airline's
    reservation data sets fares by demand.
\end{itemize}}

\Q{Data vs information vs knowledge vs wisdom; how data and information are used in an office}{\m{4}}
{\color{sub}\footnotesize \tP{3} \yr{82 Ba, \textbf{78 Bh}, 75 Ash}}
\sbs{%
\begin{itemize}
  \item \textbf{Data}: raw, unprocessed facts and figures; no meaning alone. Eg 45, 62, 38.
  \item \textbf{Information}: data processed into a meaningful, useful form. Eg ``average
    daily sales this week: Rs 48{,}000, down 12\%''.
  \item \textbf{Knowledge}: information combined with experience and understanding: knows
    \textit{how}. Eg ``sales drop every exam month''.
  \item \textbf{Wisdom}: applying knowledge with judgement to decide well: knows
    \textit{why} and \textit{what is best}. Eg ``run a student discount in exam months''.
  \item DIKW pyramid: each level adds context and value. \added
\end{itemize}
\lead{Qualities of good information}
{\small Accuracy, timeliness, completeness, conciseness, relevance, right frequency,
understandability, cost less than its value.}}{%
\lead{Use in an office (78 Bh, 75 Ash)}
\begin{itemize}
  \item \textbf{Data} recorded daily: attendance, invoices, receipts, stock movements,
    letters.
  \item Processed into \textbf{information}: payroll, monthly accounts, stock reports,
    pending-file lists.
  \item Used for: paying staff, ordering supplies, answering customers, monitoring work,
    reporting to management, planning budgets.
  \item Data $\rightarrow$ processing $\rightarrow$ information, with feedback to correct
    collection.
\end{itemize}}

\penalty0\vspace{2pt}

%=======================================================================
\T{5.2 MIS for Planning, Decision Making and Management Levels}

\Q{Role of information in the planning process and in decision making; information system is vital for planning and control}{\m{4}\m{5}\m{6}\m{8}}
{\color{sub}\footnotesize \tF{7} \yr{\textbf{76 Ch}, 76 Ash, \textbf{68 Ch}, \textbf{68 Ba}, \textbf{67 Asa}, \textbf{66 Bh}, \textbf{65 Shr}} \gap$\cdot$\gap short note ``information system for planning process'' in 68 Ba}
\sbs{%
\lead{Information system for planning}
\begin{itemize}
  \item Planning answers: \textbf{where are we? where do we want to go? how do we get there?
    when? who? how much will it cost?}
  \item MIS supplies: past performance (where we are), forecasts of demand, costs,
    market, technology (where to go), resource data (how, who, cost).
  \item Long-range (3--5 year) plan $\rightarrow$ yearly operating plan $\rightarrow$
    individual targets: MIS tracks each.
  \item Planning is continuous; MIS keeps the plan a ``living document'' by feeding actual
    results back.
\end{itemize}
}{%
\lead{Information system for decision making}
\begin{itemize}
  \item Decision = choosing the best alternative; more and better information $\rightarrow$
    better choice.
  \item Steps supported: identify problem (exception reports) $\rightarrow$ generate
    alternatives (data, models) $\rightarrow$ evaluate (what-if, DSS) $\rightarrow$ choose
    $\rightarrow$ monitor (feedback).
  \item \textbf{Strategic}: future-oriented, long range, unstructured: new plant, new market.
  \item \textbf{Tactical}: short term: budgets, fund flow, personnel problems, product
    improvement.
  \item \textbf{Operational}: structured, routine: inventory, scheduling, allocating workers.
\end{itemize}}

\sbs{%
\lead{Vital for control (68 Ch)}
\begin{itemize}
  \item Control = compare actual with plan $\rightarrow$ correct. Without timely
    information there is nothing to compare.
  \item Exception reports highlight only deviations (management by exception).
\end{itemize}}{%
\lead{Value of MIS in the planning process (76 Ch, 76 Ash)}
\begin{itemize}
  \item Accurate forecasts; realistic targets; scenario analysis; early warning when plans
    slip; coordination of departmental plans.
\end{itemize}}

\Q{Hierarchy of information needs; organizational structure and MIS; MIS at different levels of management}{\m{5}\m{8}}
{\color{sub}\footnotesize \tF{4} \yr{80 Ba, \textbf{68 Ch}, \textbf{66 Bh}, \textbf{65 Shr}} \gap$\cdot$\gap from S.K.\ Joshi, Ch.\,3.1.1}
\sbs{%
\begin{itemize}
  \item Managers at different levels need \textbf{different kinds} of information at
    \textbf{different intervals}: the \textbf{hierarchy of information needs}. MIS is a
    pyramid matching the organization's levels.
  \item \textbf{Operational (lower)}: day-to-day production of goods and services.
  \begin{itemize}
    \item Detailed, internal, precise, historical data.
    \item Hourly or daily: raw material needs, work schedules, output vs schedule.
  \end{itemize}
  \item \textbf{Tactical (middle)}: puts top's plans into operation.
  \begin{itemize}
    \item Summaries: schedules, revenue, profit, cost, other economic measures.
    \item Weekly or monthly; used to set controls and allocate resources.
  \end{itemize}
  \item \textbf{Strategic (top)}: corporate policy, growth, survival.
  \begin{itemize}
    \item Aggregated, future-oriented, much of it \textbf{external} (competitors,
      suppliers, consultants, media).
    \item Quarterly or yearly: are goals met, is the firm profitable, what should change.
  \end{itemize}
  \item Bottom of the pyramid: \textbf{transaction processing} data feeding all levels.
\end{itemize}}{%
\lead{Why the hierarchy is necessary (68 Ch)}
\begin{itemize}
  \item Each level gets \textbf{only what it needs}: no overload at top, no gaps below.
  \item Detail is condensed upward (reports), direction flows downward (plans).
  \item Matches decision type to information type.
\end{itemize}
\lead{Organizational structure and MIS}
\begin{itemize}
  \item Nature of the firm (profit or not, goods or services, single site or global)
    decides its information needs and systems.
  \item MIS draws on organization concepts: structure, communication, power, decision
    making, motivation, leadership.
\end{itemize}
{\small
\begin{tabularx}{\linewidth}{@{}L{1.4cm}XXX@{}}
\toprule
\hrow \thead{Level} & \thead{Scope} & \thead{Source} & \thead{Frequency} \\
\midrule
Strategic & Wide, summary, future & Mostly external & Quarterly, yearly \\
Tactical & Summaries, trends & Internal & Weekly, monthly \\
Operational & Narrow, detailed, current & Internal & Hourly, daily \\
\bottomrule
\end{tabularx}}}

\creamq{Information architecture}{\m{3}\m{4}}
{\color{sub}\footnotesize \tP{2} \yr{\textbf{68 Ba}, \textbf{65 Shr}} \gap$\cdot$\gap from S.K.\ Joshi, Ch.\,3.2}
\sbs{%
\begin{itemize}
  \item From Greek \textit{archos} (chief) + \textit{tekton} (carpenter): design and
    structure.
  \item Information architecture = \textbf{how the parts of an information system
    interconnect and communicate}: the general \textbf{flow} of information, how
    \textbf{databases interrelate}, and how applications make access easier for managers.
  \item Also (AJ Sir): organizing, structuring and labelling content so users find
    information and complete tasks.
\end{itemize}}{%
\begin{itemize}
  \item \textbf{Anthony's triangle}: strategic, tactical, operational applications stacked
    by level.
  \item Add \textbf{divisions} (production, marketing, accounting) $\rightarrow$ a
    two-dimensional structure. Each division keeps unique data; common events (a sale) create
    \textbf{shared data} $\rightarrow$ a \textbf{shared corporate database}.
  \item Traditional bureaucracy: information flows up, over, down; divisions rarely talk
    directly. Good architecture removes these barriers.
\end{itemize}}

\penalty0\vspace{2pt}

%=======================================================================
\T{5.3 Types of Information Systems and Functional Support}

\Q{Classification of information systems (TPS, MIS, DSS, ESS) and their support to managers; DSS vs MIS; EIS for top managers}{\m{2}\m{3}\m{4}\m{5}\m{6}}
{\color{sub}\footnotesize \tF{7} \yr{\textbf{82 Bh}, \textbf{80 Bh}, 79 Ba, \textbf{78 Bh}, 75 Ash, \textbf{71 Ch}, \textbf{69 Ch}}}

\begin{itemize}
  \item \textbf{Information system}: hardware, software, data, people and procedures that
    produce information for day-to-day, short-range and long-range activities.
\end{itemize}

{\small
\begin{tabularx}{\textwidth}{@{}L{2.4cm}XXXX@{}}
\toprule
\hrow \thead{Basis} & \thead{TPS} & \thead{MIS} & \thead{DSS} & \thead{ESS / EIS} \\
\midrule
Level served & Operational & Middle (management) & Middle, professionals & Top (strategic) \\
Inputs & Transactions, events & Summarized transaction data & Transaction data + models, external data & Aggregate internal + external data \\
Processing & Sort, list, update, batch or online & Simple models, routine reports & Interactive, what-if analysis & Interactive, graphics, drill-down \\
Outputs & Detailed reports, lists & Summary and exception reports & Decision analysis & Projections, dashboards \\
Decisions & Structured & Structured, recurring & Semi-structured & Unstructured \\
Example & Payroll, order entry, billing & Annual budgeting, sales by region & Contract cost analysis, production scheduling & 5-year operating plan \\
\bottomrule
\end{tabularx}}

\sbsr{0.58}{%
\lead{DSS vs MIS (78 Bh, 75 Ash)}
{\small
\begin{tabularx}{\linewidth}{@{}L{1.6cm}XX@{}}
\toprule
\hrow \thead{Basis} & \thead{MIS} & \thead{DSS} \\
\midrule
Purpose & Routine reporting, control & Support a specific decision \\
Problems & Structured & Semi-/unstructured \\
Output & Fixed, periodic reports & Ad hoc, interactive analysis \\
User role & Receives reports & Manipulates data and models \\
Data & Internal, past & Internal + external, future \\
Flexibility & Low & High (what-if) \\
\bottomrule
\end{tabularx}}
\lead{How production managers use DSS in addition to MIS (82 Bh)}
\begin{itemize}
  \item MIS tells \textbf{what happened}: output, scrap, machine downtime reports.
  \item DSS asks \textbf{what if}: best production schedule and product mix (linear
    programming), make-or-buy, capacity expansion, inventory reorder levels, maintenance
    timing, with simulations of cost and delay.
\end{itemize}
\lead{Significance of EIS for top managers (79 Ba)}
\begin{itemize}
  \item One-screen dashboards of key indicators; drill-down to detail.
  \item External data: competitors, markets, economy.
  \item Fast, easy to use without technical skill $\rightarrow$ strategic decisions,
    early warning, less dependence on staff reports.
\end{itemize}}{%
\figT{c5_interrel.png}
{\color{sub}\scriptsize TPS feeds MIS and DSS; all feed ESS. Adhikari notes, p299.}
\lead{Support to managers}
\begin{itemize}
  \item TPS: records every transaction; the data source for all others.
  \item MIS: regular detail, summary and exception reports $\rightarrow$ control
    (also called management reporting system).
  \item DSS: models + data to decide non-routine problems.
  \item ESS: strategic view for top managers.
  \item \textbf{Expert system}: knowledge base + inference rules imitate an expert (Eg
    medical or fault diagnosis).
  \item Modern software integrates them (ERP).
\end{itemize}}

\par\Needspace{34\baselineskip}
\Q{Information support for functional areas of management; how MIS is used at different levels}{\m{4}\m{5}\m{6}\m{8}}
{\color{sub}\footnotesize \tF{5} \yr{\textbf{74 Ch}, 73 Shr, 72 Ka, \textbf{70 Ch}, 70 Asa} \gap$\cdot$\gap 80 Ba's ``MIS at different levels'' is answered by the same figure}

\figC{c5_functional.png}{0.80\textwidth}
{\color{sub}\scriptsize Four functional areas (columns) crossed with three levels (rows): each cell
is the information system that supports it. Adhikari notes, p300.}

\sbs{%
\lead{Sales and marketing}
\begin{itemize}
  \item TPS: order processing, invoicing.
  \item MIS: sales by region, salesperson, product; sales management.
  \item DSS: pricing, sales region analysis, promotion planning.
  \item ESS: 5-year sales trend forecast, market share.
\end{itemize}
\lead{Manufacturing and production}
\begin{itemize}
  \item TPS: material movement, machine control, purchase orders.
  \item MIS: inventory control, production scheduling reports, quality and scrap.
  \item DSS: production planning, relocation cost, capacity.
  \item ESS: 5-year operating plan, new plant decisions.
\end{itemize}}{%
\lead{Finance and accounting}
\begin{itemize}
  \item TPS: payroll, accounts payable and receivable, general ledger.
  \item MIS: annual budgeting, cost reports.
  \item DSS: profitability, contract cost analysis, investment appraisal.
  \item ESS: profit planning, long-term financing.
\end{itemize}
\lead{Human resources}
\begin{itemize}
  \item TPS: employee record keeping, attendance.
  \item MIS: turnover, absenteeism, training reports, appraisal history.
  \item DSS: compensation analysis, manpower allocation.
  \item ESS: personnel (manpower) planning.
\end{itemize}}

\penalty0\vspace{2pt}

%=======================================================================
\T{5.4 Computers, Databases, Networks and Organizing Information Systems}

\Q{Computers and MIS; role of computers in MIS; what is a website}{\m{2}\m{4}}
{\color{sub}\footnotesize \tF{4} \yr{\textbf{81 Bh}, 76 Ash, 74 Ash, \textbf{68 Ba}}}
\sbs{%
\begin{itemize}
  \item Computers are the \textbf{means} of modern MIS, not MIS itself: a computer alone does
    not make work efficient; identify the management problem first.
  \item \textbf{Before computers} information was: too late, incomplete, costlier than its
    worth, unfocused, irrelevant, rigid.
  \item \textbf{Role of computers}: fast, accurate processing of huge data; storage and
    retrieval; any report format; communication across sites; models and simulation.
  \item \textbf{Quick-response systems}:
  \begin{itemize}
    \item \textbf{Online processing}: manager interacts directly with the computer.
    \item \textbf{Real-time processing}: system runs alongside the activity and guides it
      (Eg ATM, airline booking).
    \item \textbf{Time sharing}: many users share one computer at high speed.
  \end{itemize}
\end{itemize}}{%
\begin{itemize}
  \item Two uses: \textbf{data (transaction) processing} with fixed reports, and
    \textbf{end-user computing} (managers, staff using spreadsheets, queries) with
    flexibility.
  \item \textbf{Relationship}: MIS defines what information is needed; the computer
    produces it quickly and cheaply. Computer = tool, MIS = purpose.
  \item \textbf{Website} (74 Ash): a set of linked web pages under one domain name, hosted
    on a server and reached through the internet. For MIS: publishes information to
    customers, collects orders and feedback (data), and gives staff remote access to
    reports. \added
  \item \textbf{Computer-aided advertising} (69 Ch, short note \tL{1} \yr{\textbf{69 Ch}}): computer
    graphics and design for ads; targeted online, email and social media ads; analytics of
    clicks and conversions. \added
\end{itemize}}

\sbs{%
\lead{Database information system \m{4}\m{5}}
{\color{sub}\footnotesize \tF{4} \yr{81 Ba, \textbf{68 Ba}, \textbf{66 Bh}, \textbf{65 Shr}} \gap$\cdot$\gap short note in 68 Ba}
\begin{itemize}
  \item \textbf{Database}: organized collection of related data serving many applications at
    once, appearing to be in one place.
  \item \textbf{Personal database}: an individual's contacts, schedules, notes.
    \textbf{Corporate database}: large, shared, integrity critical.
  \item \textbf{DBMS}: software to create, maintain and query it (Oracle, MySQL); applications
    extract data without their own files.
  \item Components: data dictionary (description), architecture (relations), the objects
    described, rules for manipulation.
  \item $+$ no redundancy, consistency, sharing, security, easy reporting.
  \item Eg a bank's customer database used by accounts, loans, ATM and mobile banking;
    a college's student database used by admission, exams, library, fees.
\end{itemize}}{%
\lead{Networking information system \m{4}\m{5}}
{\color{sub}\footnotesize \tP{2} \yr{\textbf{67 Asa}, \textbf{65 Shr}} \gap$\cdot$\gap short note in 65 Shr}
\begin{itemize}
  \item Telecommunication network linking computers and terminals so users at different
    locations share data and resources.
  \item Components: computers, terminals, channels, communication processors, software;
    common rules = \textbf{protocols}; analog or digital signals.
  \item Types: LAN, MAN, WAN, internet, intranet, extranet.
  \item Needs a flexible, standard network architecture that grows with the business.
  \item $+$ real-time sharing, distributed offices, e-commerce, email, lower cost.
  \item Eg bank branches and ATMs on one network; NEA billing linked across offices.
\end{itemize}}

\Q{How information systems can be organized properly}{\m{4}\m{6}}
{\color{sub}\footnotesize \tP{3} \yr{\textbf{81 Bh}, 76 Ash, 71 Shr} \gap$\cdot$\gap short note in 71 Shr}
\sbs{%
\begin{itemize}
  \item \textbf{Complete centralization}: one central IS department holds hardware,
    software, data and MIS personnel; serves all user departments.
  \begin{itemize}
    \item $+$ control, standards, economy of scale, security.
    \item $-$ slow response, users feel distant, bottleneck.
  \end{itemize}
  \item \textbf{Complete decentralization}: each department has its own systems and MIS staff.
  \begin{itemize}
    \item $+$ responsive, users own their systems.
    \item $-$ duplication, incompatible data, higher cost.
  \end{itemize}
  \item \textbf{Hybrid (distributed)}: central database and standards, local processing and
    staff. Most common today.
  \item Also organize by: steering committee of managers, CIO at top level, clear
    policies for data ownership and security, user training.
\end{itemize}}{%
\figT{c5_central.png}
{\color{sub}\scriptsize Complete centralization. Adhikari notes, p301.}
\vspace{2pt}
\figT{c5_decentral.png}
{\color{sub}\scriptsize Complete decentralization: MIS staff sit in user departments. Adhikari
notes, p302.}}

\creamq{Evolution of MIS}{}
{\color{sub}\footnotesize \tN}
\begin{itemize}
  \item Manual ledgers and punch cards $\rightarrow$ electronic data processing (1950s--60s,
    record keeping) $\rightarrow$ MIS reports (1960s--70s) $\rightarrow$ DSS (1970s) $\rightarrow$
    EIS / ESS (1980s) $\rightarrow$ PCs and end-user computing $\rightarrow$ networks, ERP,
    internet, cloud and business intelligence.
  \item Driven by: growth of management theory, management science, new production methods
    and structures, and computers.
\end{itemize}

\penalty0\vspace{2pt}

%=======================================================================
\T{5.5 PYQ Mapping}

\lead{Theory questions, covered in this document}
\begin{itemize}
  \item \textbf{Define MIS; need, functions, importance; MIS beyond task software}
    $\rightarrow$ $\S$5.1. 17 papers; plus the second half of 73 Ch and 71 Ch, which the
    Detailed PYQ files under case studies.
  \item \textbf{Data, information, knowledge, wisdom; use in an office} \yr{82 Ba \m{4},
    \textbf{78 Bh} \m{4}, 75 Ash \m{4}} $\rightarrow$ $\S$5.1.
  \item \textbf{Information for planning, decision making, control} $\rightarrow$ $\S$5.2.
    7 papers.
  \item \textbf{Hierarchy of information needs; MIS at levels} \yr{80 Ba \m{8}, \textbf{68 Ch} \m{5},
    \textbf{66 Bh} \m{5}, \textbf{65 Shr} \m{4}}; \textbf{information architecture} \yr{\textbf{68 Ba} \m{3},
    \textbf{65 Shr} \m{4}} $\rightarrow$ $\S$5.2.
  \item \textbf{Types of IS, DSS vs MIS, DSS for production, EIS} $\rightarrow$ $\S$5.3.
    7 papers.
  \item \textbf{Information support for functional areas} $\rightarrow$ $\S$5.3.
    \yr{\textbf{74 Ch} \m{8}, 73 Shr \m{4}, 72 Ka \m{8}, \textbf{70 Ch} \m{6}, 70 Asa \m{5}}
  \item \textbf{Computers and MIS; website; computer-aided advertising} $\rightarrow$
    $\S$5.4. \yr{\textbf{81 Bh}, 76 Ash, 74 Ash, \textbf{69 Ch}, \textbf{68 Ba}}
  \item \textbf{Database IS} \yr{81 Ba \m{5}, \textbf{68 Ba} \m{5}, \textbf{66 Bh} \m{5}, \textbf{65 Shr} \m{4}};
    \textbf{networking IS} \yr{\textbf{67 Asa} \m{5}, \textbf{65 Shr} \m{4}}; \textbf{organizing IS}
    \yr{\textbf{81 Bh} \m{6}, 76 Ash \m{6}, 71 Shr \m{4}} $\rightarrow$ $\S$5.4.
  \item \textbf{Evolution of MIS}: syllabus, never set $\rightarrow$ $\S$5.4.
\end{itemize}

\lead{Numerical content}
\begin{itemize}
  \item \textbf{This chapter has no numerical component.}
\end{itemize}

\lead{Sources and corrections}
\begin{itemize}
  \item Hierarchy of information needs, information architecture, planning and the
    quick-response systems come from S.K.\ Joshi, Ch.\,3 (pp\,99--135), read through
    \code{tools/ocr\_page.py}; the lecture notes do not cover them.
  \item Adhikari p283 expands DSS as ``design support system''; it is \textbf{decision}
    support system. Adhikari p296 heads the ESS slide but describes an \textbf{expert}
    system; both are given separately above.
  \item The Detailed PYQ's ``hierarchy of needs'' in 65 Shr sits in the MIS question and is
    read as hierarchy of information needs, as 66 Bh words it.
\end{itemize}
